% ---------------------------------------------------------------------------------------
% MODELO DE TCC - INSTITUTO FEDERAL DE SERGIPE (IFS)
% ---------------------------------------------------------------------------------------
% Autor do Modelo: Mikael dos Santos Gois (mikaelgois1@gmail.com)
% Baseado nos modelos de: 
%   - Diogo D. Moreira (TCC, Dissertação e Tese)
%   - DECOM/UFOP (Monografia de Ciência da Computação)
% Modelo Base: abnTeX2 adaptado para as normativas internas de TCC do IFS (Art. 31 e 34)
% Versão: 1.0 (2026)
% ---------------------------------------------------------------------------------------
% Este arquivo faz parte de um modelo opensource para estudantes do IFS.
% Não remova os créditos dos autores originais e do modelo base.
% Este modelo é fornecido "no estado em que se encontra", sem garantias de qualquer tipo.
% ---------------------------------------------------------------------------------------

\begin{resumo}[Abstract]
  \begin{otherlanguage*}{english}
    Programming is a fundamental skill in the 21st century, but introductory teaching faces significant
    barriers that discourage initial learning, such as the English language barrier and premature syntactic
    complexity. This work presents the Luso project, an educational programming language designed with
    syntax in Portuguese. Based on the theories of Software Localization and Gradual Programming, Luso
    aims to reduce the cognitive load of novice students, serving as a scaffolding for the future transition
    to professional languages such as Java and Python. The project proposes an ecosystem that includes
    a transpiler, an IDE, and support for third-party development tools, seeking to align technology
    with the teacher's teaching methods. In addition, the project has educational documentation for users
    and technical documentation for project contributors.

    \vspace{\onelineskip}

    \noindent
    \textbf{Keywords}: Teaching Programming; Gradual Languages; Software Localization; Syntax Barrier; Cognitive Load.
  \end{otherlanguage*}
\end{resumo}